\par\bigskip\pdfbookmark[0]{Über den Zusammenhang. Fragment VI}{nexus-fragment-VI}\centerline{[VI.]}\medskip
\textsc{Satz.} Die Anzahl der Klassen, in die alle eigentlich primitiven binären Formen mit der negativen Determinante $-D$\begingroup\renewcommand{\thefootnote}{*)}\footnote{[eingeteilt werden.]}\endgroup, ist gleich
\[
\frac\pi4\times\surd D\times\text{Produkt aus }\frac{p-1}p\frac{q+1}q\times\frac{rr-1}{rr}
\]
wobei bezeichnet:
\begin{description}
\item[$p$] alle Primzahlen\begingroup\renewcommand{\thefootnote}{**)}\footnote{[ungerade.]}\endgroup, für die $-D$ ein Nichtrest ist,
\item[$q$] alle Primzahlen\textsuperscript{**)}, für die $-D$ ein Rest ist,
\item[$r$] alle Primzahlen\textsuperscript{**)}, die $D$ teilen;
\end{description}
\[
=\frac{\frac\pi4\surd D\ \text{Produkt aus }\frac{rr-1}{rr}}{1\pm\tfrac13\pm\tfrac15\text{ usw.}}
\]

Im Nenner steht das positive Vorzeichen vor den Brüchen, deren Nenner die Form von Nichtteilern haben, das negative vor denen, deren Nenner die Form von Teilern von $xx+D$ haben; die Brüche, deren Nenner nicht zu $D$ teilerfremd sind, werden ganz weggelassen\begingroup\renewcommand{\thefootnote}{***)}\footnote{[Bezeichnet man mit $m$ alle positiven ganzen Zahlen, die zu $2D$ teilerfremd sind, und verwendet das von \textsc{Jacobi} verallgemeinerte Symbol von \textsc{Legendre}, so ist die obige Regel für die Vorzeichen folgendermaßen zu berichtigen: In der vorhergehenden Formel ist der Nenner
\[
1\pm\tfrac13\pm\tfrac15\text{ usw.}=\Sigma\pm\left(\frac{-D}m\right)\frac1m
\]
wobei das obere oder untere Vorzeichen zu nehmen ist, je nachdem die Zahl $m$ ein Produkt aus einer geraden oder ungeraden Anzahl (gleicher oder ungleicher) Primzahlen ist; dagegen ist im Zähler der nachfolgenden Formel
\[
1\pm\tfrac13\pm\tfrac15\ldots=\Sigma\left(\frac{-D}m\right)\frac1m
\]
]}\endgroup.

\SourceLeaf{296}
\[
\begin{aligned}
&=\frac{2\surd D(1\pm\tfrac13\pm\tfrac15\ldots)}\pi\\
&=\frac{\operatorname{cotg}\theta\pm\operatorname{cotg}3\theta\pm\operatorname{cotg}5\theta\ldots\pm\operatorname{cotg}n\theta}{N:\surd D}
\end{aligned}
\]
wobei $\theta=\frac\pi N$, $N=\left\{\begin{smallmatrix}1\\4\end{smallmatrix}\right\}D$ gesetzt und für $n$ alle zu $D$ teilerfremden Zahlen genommen werden, mit dem wie oben bestimmten Vorzeichen\begingroup\renewcommand{\thefootnote}{*)}\footnote{[Siehe die weiter unten folgende Anmerkung zu diesem Fragment.]}\endgroup.

Für eine positive Determinante ist die Anzahl der Klassen\begingroup\renewcommand{\thefootnote}{**)}\footnote{[In der nachfolgenden Formel bezeichnet $D$ die positive Determinante, und es ist
\[
1\pm\tfrac13\pm\tfrac15\ldots=\Sigma\left(\frac Dm\right)\frac1m
\]
]}\endgroup
\[
=\frac{2\surd D(1\pm\tfrac13\pm\tfrac15\ldots)}{\log .T+U\surd D}
\]
Dabei bezeichnen $T,U$ die kleinsten Werte der Größen $t,u$, die der Gleichung $tt-Duu=1$ genügen
\[
=\frac{\log\sin\tfrac12\theta\pm\log\sin\tfrac32\theta\pm\log\sin\tfrac52\theta\text{ usw.}}{\log T+U\surd D}
\]
\BookRule{.30\textwidth}
